\documentclass[11pt]{amsart}
\ifdefined\XeTeXrevision\else\ifdefined\pdfoutput\pdfoutput=1\fi\fi

\usepackage[margin=1.2in]{geometry}
\usepackage{amsmath,amssymb,mathtools}
\usepackage{booktabs}
\usepackage{xcolor}
\definecolor{linkblue}{RGB}{25,60,120}
\usepackage[colorlinks=true,linkcolor=linkblue,citecolor=linkblue,urlcolor=linkblue]{hyperref}
\hypersetup{
  pdftitle={No triangle can be cut into nineteen congruent triangles: the prime case of Erdos Problem 634},
  pdfauthor={Jan Philipp Harries},
  pdfsubject={Triangle tilings by congruent triangles},
  pdfkeywords={triangle tilings, congruent triangles, prime tilings, Erdos problem 634}
}

\newtheorem{theorem}{Theorem}
\newtheorem{lemma}[theorem]{Lemma}
\newtheorem{corollary}[theorem]{Corollary}
\theoremstyle{remark}
\newtheorem{remark}[theorem]{Remark}

\newcommand{\Z}{\mathbb{Z}}

\title[The prime case of Erd\H{o}s Problem 634]{No triangle can be cut into nineteen congruent triangles:\\ the prime case of Erd\H{o}s Problem 634}

\author{Jan Philipp Harries}
\thanks{ellamind GmbH, Bremen, Germany. \textit{E-mail address:} \href{mailto:jan@ellamind.com}{jan@ellamind.com}}

\date{July 24, 2026}

\subjclass[2020]{Primary 52C20; Secondary 51M04, 11D09}
\keywords{triangle tilings, congruent triangles, Erd\H{o}s problem 634, prime tilings}

\begin{document}

\begin{abstract}
We prove that for every prime $p>3$ with $p\equiv 3 \pmod 4$, no nondegenerate Euclidean triangle can be tiled by exactly $p$ pairwise interior-disjoint congruent triangles. In particular, no triangle can be cut into nineteen congruent triangles, answering a question highlighted in Erd\H{o}s Problem~634. Combined with the classical constructions realizing $p=2$, $p=3$, and every prime $p\equiv 1\pmod 4$, this determines the prime tile counts completely: a triangle can be tiled by a prime number $p$ of congruent triangles if and only if $p=2$, $p=3$, or $p\equiv 1\pmod 4$. The proof combines recent classification theorems of Beeson, Laczkovich, and Zhang with an elementary arithmetic argument excluding a prime tile count in the remaining configurations, in which the tile has a $2\pi/3$ angle. Several of the structural results we use are at present available only as preprints.
\end{abstract}

\maketitle

\section{Introduction}

Any triangle can be cut into $n^2$ congruent triangles similar to itself, for every $n\ge 1$. Erd\H{o}s asked, in a question reported by Soifer \cite{Soifer09}, for which $N$ there exists \emph{some} triangle that can be cut into $N$ congruent triangles; this is Problem~634 in Bloom's collection \cite{EP634}. Soifer showed that all $N$ of the forms $n^2$, $2n^2$, $3n^2$, $6n^2$, and $n^2+m^2$ are realizable, and Zhang \cite{Zhang25} recently gave further constructive families. In the other direction, Beeson proved that no triangle can be cut into $7$ or into $11$ congruent triangles \cite{Beeson7}, and in July 2026 that no \emph{isosceles} triangle can be cut into a prime number $N>3$ of congruent triangles \cite{BeesonPrime}. The page \cite{EP634} states: ``It is possible that any prime of the form $4n+3$ does not have this property. In particular, it is not known if $19$ has this property.''

We settle this question.\footnote{A withdrawn preprint of Beeson \cite{Beeson2228} had claimed, among other results, that no $19$-tiling exists. It was withdrawn in 2024 after its main theorem was found to be false; the arXiv record notes that tilings described in later papers contradict it. The present proof makes no use of that manuscript.}

\begin{theorem}\label{thm:main}
Let $p>3$ be a prime with $p\equiv 3\pmod 4$. No nondegenerate Euclidean triangle can be tiled by exactly $p$ pairwise interior-disjoint congruent triangles. In particular, no triangle can be cut into $19$ congruent triangles.
\end{theorem}

Here a \emph{tiling} of a triangle $T$ by a triangle $R$ is a finite family of triangles congruent to $R$ (reflections allowed), with pairwise disjoint interiors, whose union is $T$; the tiling is not assumed edge-to-edge. Following \cite{BLZ633}, a tiling is a \emph{reptiling} if $R$ is similar to $T$; a triangle has \emph{commensurable angles} if all of its angles are rational multiples of $\pi$, and \emph{incommensurable angles} otherwise; and a triangle has \emph{commensurable sides} if the pairwise ratios of its side lengths are rational. We call a triangle \emph{scalene} if no two of its sides are equal; ``isosceles'' is used inclusively, so that equilateral triangles are isosceles and ``not isosceles'' means scalene.

Since every prime $p\equiv 1\pmod 4$ is a sum of two squares (Fermat), the classical hypotenuse construction of Snover, Waiveris, and Williams \cite{SWW91} realizes such $p$ as tile counts for a right triangle, and $2$ and $3$ are realized by bisecting an isosceles right triangle and by the $3$-tiling of the equilateral triangle. Theorem~\ref{thm:main} is therefore sharp and yields:

\begin{corollary}\label{cor:primes}
A prime $p$ is the tile count of some tiling of some triangle by congruent triangles if and only if $p=2$, $p=3$, or $p\equiv 1\pmod 4$.
\end{corollary}

This completes the prime case of Erd\H{o}s Problem~634. In passing it recovers the cases $N=7$ and $N=11$, first proved by Beeson \cite{Beeson7} by different methods.

The proof of Theorem~\ref{thm:main} is short because it stands on recent structural results of Beeson, Laczkovich, and Zhang, recalled in Section~\ref{sec:imports}, which reduce the problem to tiles with a $2\pi/3$ angle and commensurable sides. The remaining arithmetic input is Lemma~\ref{lem:group2} in Section~\ref{sec:lemma}. We do not claim novelty for its algebraic ingredients: the side vectors and tile-count expressions occurring there appear in Zhang's constructions \cite{Zhang25}, and the basic arithmetic of the equation $c^2=a^2+ab+b^2$ is developed by Beeson in \cite{BeesonPrime}; detailed attributions are given in Section~\ref{sec:lemma}. The point of Lemma~\ref{lem:group2} is the necessity statement needed here: in each of the four remaining configurations, for an \emph{arbitrary} tiling, boundary integrality and primitivity force the scale factor between tile and tiled triangle to be an integer, and the resulting tile counts are never prime. Section~\ref{sec:proof} assembles the proof, and Section~\ref{sec:verification} describes supplementary computational checks of the algebra.

Theorem~\ref{thm:main} relies on several recent results that are currently available only as preprints---\cite{BLZ633,BZRat,BeesonPrime} and, through Theorem~\ref{thm:group1prime} below, the manuscript \cite{Beeson2229}---and these have not all appeared in refereed venues at the time of writing. The precise statements used are reproduced in Section~\ref{sec:imports}; we use them as black boxes and do not re-derive their hypotheses or conclusions.

\section{Imported results}\label{sec:imports}

Throughout, $T$ denotes the tiled triangle and $R$ the tile. We use the following results.

\begin{theorem}[{Laczkovich \cite{Lacz95,Lacz12}; stated as \cite[Theorem~3]{BLZ633}}]\label{thm:commensurable}
Suppose $T$ has commensurable angles and is tiled by $R$. If $T$ is not an isosceles (or equilateral) triangle, then the tiling is a reptiling.
\end{theorem}

\begin{theorem}[{Snover--Waiveris--Williams \cite{SWW91}; stated as \cite[Theorem~4]{BLZ633}}]\label{thm:reptile}
If $T$ can be reptiled into $N$ tiles, then $N=M^2$, or $N=M^2+K^2$, or $N=3M^2$ for some positive integers $M,K$.
\end{theorem}

\begin{theorem}[{Beeson--Laczkovich--Zhang \cite[Theorem~9]{BLZ633}, from Laczkovich \cite[Theorem~4.1]{Lacz95}}]\label{thm:classification}
Let a non-isosceles triangle $T$ be tiled by a tile $R$ with incommensurable angles $(\alpha,\beta,\gamma)$. Then one of the following holds:
\begin{enumerate}
\item the tiling is a reptiling;
\item \textup{(Group 1)} $3\alpha+2\beta=\pi$, and $T$ has angles $(\alpha,2\alpha,2\beta)$ or $(2\alpha,\beta,\alpha+\beta)$;
\item \textup{(Group 2)} $\gamma=2\pi/3$, and $T$ has angles
\[
(\alpha,2\alpha,3\beta),\quad
(\alpha,2\beta,2\alpha+\beta),\quad
(\alpha,\alpha+\beta,\alpha+2\beta),\quad\text{or}\quad
(2\alpha,2\beta,\alpha+\beta).
\]
\end{enumerate}
\end{theorem}

\begin{theorem}[{Beeson \cite[Theorem~1]{BeesonPrime}, proved in \cite{Beeson2229}}]\label{thm:group1prime}
Let $T$ be tiled by $N$ copies of a tile $R$ with angles $(\alpha,\beta,\gamma)$ satisfying $3\alpha+2\beta=\pi$, where $T$ is not similar to $R$. Then $N$ is not prime.
\end{theorem}

\begin{theorem}[{Beeson--Zhang \cite[Theorem~1.2]{BZRat}; stated as \cite[Theorem~6]{BLZ633}}]\label{thm:rationality}
Let $T$ be tiled by a tile $R$ such that $R$ is not similar to $T$, $R$ is not a right triangle, and $R$ has incommensurable angles. Then $R$ has commensurable sides.
\end{theorem}

\begin{theorem}[{Beeson \cite[Theorem~11]{BeesonPrime}}]\label{thm:isosceles}
Let $T$ be an isosceles triangle tiled into $N$ copies of a triangle $R$. Then $N$ is not a prime greater than three.
\end{theorem}

In Theorem~\ref{thm:classification}, the labeling of the two acute tile angles $\alpha,\beta$ in Group~2 is not ordered, so the four listed triples are to be read up to the interchange $\alpha\leftrightarrow\beta$. Lemma~\ref{lem:group2} below is symmetric under this interchange, and the computational checks of Section~\ref{sec:verification} run both labelings.

\section{An arithmetic lemma for tiles with a \texorpdfstring{$2\pi/3$}{2pi/3} angle}\label{sec:lemma}

\begin{lemma}\label{lem:group2}
Let $R$ have angles $(\alpha,\beta,\gamma)$ with $\gamma=2\pi/3$ and commensurable sides, scaled so that the sides opposite $(\alpha,\beta,\gamma)$ are positive integers $(a,b,c)$ with $\gcd(a,b,c)=1$; we call such a triple \emph{primitive}. If $T$ has one of the four Group~2 angle triples of Theorem~\ref{thm:classification} and is tiled by $N$ copies of $R$, then $N$ is not prime.
\end{lemma}

\begin{proof}
The law of cosines with $\gamma=2\pi/3$ gives
\begin{equation}\label{eq:norm}
c^2=a^2+ab+b^2 .
\end{equation}
Primitivity implies pairwise coprimality: a prime dividing two of $a,b,c$ divides the third by \eqref{eq:norm} (cf.\ \cite[Lemma~5]{BeesonPrime}). Moreover
\begin{equation}\label{eq:not3}
3\nmid c :
\end{equation}
if $3\mid c$, then \eqref{eq:norm} modulo $3$ gives $(a-b)^2\equiv 0\pmod 3$, so $a=b+3k$; substituting, $c^2=3b^2+9bk+9k^2$, and $9\mid c^2$ forces $3\mid b$ and then $3\mid a$, contradicting primitivity (cf.\ the proof of \cite[Theorem~6]{BeesonPrime}).

\subsection*{Boundary integrality}
With this scaling, every side of $T$ has integer length; no edge-to-edge assumption is needed. Let $S$ be a side of $T$ and let $\ell$ be the line containing $S$. Every tile is contained in $T$, hence in the closed half-plane bounded by $\ell$ that contains $T$; thus $\ell$ supports each tile it meets. If a tile meets the relative interior of $S$ in more than one point, its intersection with $\ell$ is a face of the tile---a vertex or an entire edge---and, containing two points, it is an entire edge, of length $a$, $b$, or $c$. Two such edges have disjoint relative interiors: if the edges of two tiles overlapped in a nondegenerate segment of $S$, the two tiles would occupy a common half-neighborhood of an interior point of the overlap on the inner side of $\ell$, contradicting disjoint interiors. The tiling is finite and covers $S$, and only finitely many points of $S$ lie solely in tiles meeting $S$ in a single point; hence $S$ is partitioned, up to finitely many points, into complete tile edges, and $|S|$ is a sum of elements of $\{a,b,c\}$---an integer.

\subsection*{Side and area computations}
Let $K=2c/\sqrt 3$, the circumdiameter of $R$. The laws of sines and cosines give
\[
K\sin\alpha=a,\qquad K\sin\beta=b,\qquad K\sin\gamma=c,\qquad
\cos\alpha=\frac{a+2b}{2c},\qquad \cos\beta=\frac{2a+b}{2c}.
\]
Since $\alpha+\beta=\pi/3$, the addition formulas yield
\begin{equation}\label{eq:trig}
\begin{aligned}
K\sin(\alpha+\beta)&=c, &
K\sin 2\alpha&=\frac{a(a+2b)}{c}, &
K\sin 2\beta&=\frac{b(2a+b)}{c},\\
K\sin(\alpha+2\beta)&=K\sin(2\alpha+\beta)=a+b, &
K\sin 3\beta&=\frac{3ab(a+b)}{c^{2}}. &&
\end{aligned}
\end{equation}
For a triangle with angles $(A,B,C)$ and sine-law scale $K_T$ (so its sides are $K_T\sin A$, etc.), the area equals $\tfrac{1}{2}K_T^{2}\sin A\sin B\sin C$. Writing $m>0$ for the scale factor from the displayed integer vector to the actual sides of $T$, formulas \eqref{eq:trig} give, for the four Group~2 cases:

\begin{center}
\begin{tabular}{@{}clll@{}}
\toprule
Case & Angles of $T$ & Integer side vector of $T$ & Tile count $N$ \\
\midrule
A & $(\alpha,2\alpha,3\beta)$ & $\bigl(c^2,\;c(a+2b),\;3b(a+b)\bigr)$ & $3m^2(a+2b)(a+b)$ \\
B & $(\alpha,2\beta,2\alpha+\beta)$ & $\bigl(ac,\;b(2a+b),\;c(a+b)\bigr)$ & $m^2(2a+b)(a+b)$ \\
C & $(\alpha,\alpha+\beta,\alpha+2\beta)$ & $\bigl(a,\;c,\;a+b\bigr)$ & $m^2(a+b)/b$ \\
D & $(2\alpha,2\beta,\alpha+\beta)$ & $\bigl(a(a+2b),\;b(2a+b),\;c^2\bigr)$ & $m^2(a+2b)(2a+b)$ \\
\bottomrule
\end{tabular}
\end{center}

\smallskip
Here $N$ is the ratio of areas, computed from the sine products and the ratio of sine-law scales; the scale ratios $K_T/K$ are $mc^2/a$, $mc$, $m$, and $mc$ in cases A--D respectively. Primitivity of the four vectors is proved below.

Side vectors and tile counts equivalent to the rows of this table appear in Zhang's paper \cite{Zhang25}, where they arise from explicit constructions: for case~D in \cite[Section~4]{Zhang25} (Proposition~6), and for cases B, C, and A in \cite[Propositions~10, 9, and~11]{Zhang25}, respectively. Moreover, for case~D, Zhang's Lemma~5 already proves the necessity statement we are after, under the hypothesis $a\not\equiv b\pmod 3$---which holds for every primitive triple, since $c^2\equiv(a-b)^2\pmod 3$ and $3\nmid c$ by \eqref{eq:not3}. What the present lemma adds is the same necessity statement in all four cases, together with the exclusion of a prime count.

\subsection*{The vectors are primitive}
In case C, $\gcd(a,c)=1$. In case B, a common prime divisor $q$ of all three entries satisfies $q\mid c$ (if $q\nmid c$ then $q\mid a$ and $q\mid a+b$, so $q\mid b$, impossible); then $q\nmid b$, so $q\mid 2a+b$, and reducing \eqref{eq:norm} modulo $q$ using $b\equiv-2a$ gives $3a^2\equiv 0\pmod q$, so $q=3$, contradicting \eqref{eq:not3}. In case D, a common prime $q$ divides $c^2$ but neither $a$ nor $b$, hence divides $a+2b$ and $2a+b$, hence $2(a+2b)-(2a+b)=3b$; since $q\nmid b$, this forces $q=3$, again contradicting \eqref{eq:not3}. In case A, a common prime $q$ divides $c$, hence not $3b$ by \eqref{eq:not3} and coprimality, hence $q \mid a+b$; reducing \eqref{eq:norm} modulo $q$ with $a\equiv-b$ gives $b^2\equiv 0\pmod q$, contradicting $\gcd(b,c)=1$.

\subsection*{The scale factor is an integer}
The sides of $T$ are integers and equal $m$ times the primitive vector $v=(v_1,v_2,v_3)$ of the relevant row. By B\'ezout, there are integers $r_1,r_2,r_3$ with $r_1v_1+r_2v_2+r_3v_3=1$. Each $mv_i$ is an integer, being a side length of $T$; therefore
\[
m=r_1(mv_1)+r_2(mv_2)+r_3(mv_3)\in\Z,
\]
and $m>0$ gives $m\in\Z_{>0}$.

\subsection*{Conclusion}
Since $a,b\ge 1$, we have $a+b\ge 2$, $a+2b\ge 3$, and $2a+b\ge 3$, while $m\ge 1$. In cases A, B, and D the count $N$ is therefore a product of two integer factors each greater than one---for instance $N=\bigl(m^2(2a+b)\bigr)\cdot(a+b)$ in case B---hence composite (in case A, $N$ is moreover divisible by~$3$).

In case C, suppose $N=p$ is prime. Then $pb=m^2(a+b)$. Since $\gcd(b,a+b)=\gcd(a,b)=1$, we get $b\mid m^2$; write $m^2=br$ with $r\in\Z_{>0}$. Then $p=r(a+b)$, and since $p$ is prime and $a+b\ge 2$, necessarily $r=1$ and
\[
p=a+b .
\]
Now \eqref{eq:norm} gives $c^2=(a+b)^2-ab=p^2-ab$, so $0<c<p$. Modulo $p$ we have $b\equiv-a$, hence $c^2\equiv -ab\equiv a^2 \pmod p$, so $p\mid (c-a)(c+a)$. Since $0<a,c<p$, either $c=a$ or $c+a=p$, i.e.\ $c=b$. But $c=a$ in \eqref{eq:norm} forces $b(a+b)=0$ and $c=b$ forces $a(a+b)=0$, both impossible. Hence $N$ is not prime in case C either.

All four cases (and, by the $a\leftrightarrow b$ symmetry of the argument, their relabelings) are settled.
\end{proof}

\section{Proof of Theorem~\ref{thm:main}}\label{sec:proof}

\begin{proof}
Let $p>3$ be prime with $p\equiv 3\pmod 4$ and suppose a triangle $T$ is tiled by $p$ congruent copies of a tile $R$.

\emph{Step 1: $T$ is not isosceles.} This is Theorem~\ref{thm:isosceles}, since $p>3$. In particular $T$ is scalene.

\emph{Step 2: the tiling is not a reptiling.} Otherwise Theorem~\ref{thm:reptile} gives $p=M^2$, $p=M^2+K^2$, or $p=3M^2$. A prime is not a square. If $p=3M^2$, then $3\mid p$, so $p=3$ and $M=1$, contrary to $p>3$. And $p\equiv 3\pmod 4$ is not a sum of two squares, since squares are $0$ or $1$ modulo $4$. In particular $R$ is not similar to $T$.

\emph{Step 3: $R$ has incommensurable angles.} Each corner of $T$ has angle less than $\pi$. The corner point cannot lie in the interior of a tile, where the tile would contribute angle $2\pi$, nor in the relative interior of a tile edge, where it would contribute angle $\pi$; hence it is a vertex of every tile containing it, and each angle of $T$ is a finite sum of tile angles. If all angles of $R$ were rational multiples of $\pi$, the same would hold for $T$; then Theorem~\ref{thm:commensurable} ($T$ being scalene) would force a reptiling, contradicting Step~2.

\emph{Step 4: the two remaining groups.} By Theorem~\ref{thm:classification} and Step~2, the tiling belongs to Group~1 or Group~2. In Group~1, Theorem~\ref{thm:group1prime} applies ($T$ is not similar to $R$ by Step~2) and shows $N=p$ is not prime, a contradiction. In Group~2, the tile satisfies $\gamma=2\pi/3$, so $\alpha+\beta=\pi/3$ and $0<\alpha,\beta<\pi/3$; in particular $R$ is not a right triangle. Since $R$ is not similar to $T$ and has incommensurable angles, Theorem~\ref{thm:rationality} shows $R$ has commensurable sides, and Lemma~\ref{lem:group2} shows $N=p$ is not prime, again a contradiction.

All cases are impossible, so no such tiling exists. Taking $p=19$ gives the titular statement.
\end{proof}

\section{Computational checks}\label{sec:verification}

The proofs of all assertions in this note are contained in the text; the computations described here are supplementary checks. An accompanying Python script uses exact symbolic arithmetic (SymPy) throughout, with no floating-point computation; see the ancillary files \texttt{verify\_triangle19.py} and \texttt{README.txt}.

\emph{Exact symbolic checks.} The script verifies the six trigonometric identities \eqref{eq:trig} and the four area-ratio formulas of the table symbolically, modulo the relation \eqref{eq:norm} and in both labelings of the acute angles.

\emph{Finite stress tests.} The primitivity (gcd) assertions for the four side vectors are universal statements, proved above; the script additionally tests them on $41{,}112$ primitive solutions of \eqref{eq:norm}, generated in both labelings from the two-branch parametrization of \cite[Theorem~6]{BeesonPrime} over $2\le s\le 300$, $1\le t<s$, $\gcd(s,t)=1$, $s\not\equiv t\pmod 3$. Such finite testing serves as error detection and is not a proof. As a further audit of case C for $p=19$, the script checks that all four candidate pairs $(a,b)=(19-m^2,m^2)$, $1\le m\le 4$, fail to make $c^2=a^2+ab+b^2$ a perfect square. It also verifies the identity $a+b=s(s+2t)$ on the parametrization, which yields a second route to the case C conclusion, since $s\ge 2$ and $s+2t\ge 4$ make $a+b$ composite.

\section*{Use of AI tools}

This note was produced with substantial use of AI systems, which we disclose in the interest of transparency. An initial version of Lemma~\ref{lem:group2}, of the assembly in Section~\ref{sec:proof}, and of the verification script was generated by OpenAI GPT-5.6~Pro on July~24, 2026. Claude (Anthropic, Fable~5) was used as a second-model cross-check: it re-derived the identities and arguments of Lemma~\ref{lem:group2} and Section~\ref{sec:proof}, compared the statements of the imported theorems against the cited arXiv versions, and audited and executed the verification script. Cross-checking by a second model is a useful error filter, but it is not independent verification in the scholarly sense, and neither the model-based checks nor the computations of Section~\ref{sec:verification} replace the proofs given in the text. The author reviewed the manuscript and takes full responsibility for its content. The imported Theorems~\ref{thm:commensurable}--\ref{thm:isosceles} are human mathematics due to the cited authors; no claim of AI contribution is made for them.

\section*{Acknowledgements}

This note builds directly on the program of Michael Beeson, Mikl\'os Laczkovich, and Yan~X.~Zhang, whose structural theorems carry the bulk of the mathematical weight, and on Thomas Bloom's curation of the Erd\H{o}s problems collection.

\begin{thebibliography}{99}

\bibitem{Beeson7}
M.~Beeson,
\emph{No triangle can be cut into seven congruent triangles},
arXiv:1811.09723v5 (2019).

\bibitem{Beeson2228}
M.~Beeson,
\emph{Triangle tiling V: tilings by a tile with integer sides},
arXiv:1206.2228 (2012); withdrawn by the author in 2024.

\bibitem{Beeson2229}
M.~Beeson,
\emph{Triangle tiling: the case $3\alpha+2\beta=\pi$},
arXiv:1206.2229v3 (2019).

\bibitem{BeesonPrime}
M.~Beeson,
\emph{No prime tiling of an isosceles triangle},
arXiv:2607.19572v1 (2026).

\bibitem{BLZ633}
M.~Beeson, M.~Laczkovich, and Y.~X.~Zhang,
\emph{Solution of Erd\H{o}s Problem 633},
arXiv:2604.03609v2 (2026).

\bibitem{BZRat}
M.~Beeson and Y.~X.~Zhang,
\emph{Rationality of certain triangle tilings},
arXiv:2604.01314v1 (2026).

\bibitem{EP634}
T.~F.~Bloom,
\emph{Erd\H{o}s Problem \#634},
\url{https://www.erdosproblems.com/634}, accessed July 24, 2026.

\bibitem{Lacz95}
M.~Laczkovich,
\emph{Tilings of triangles},
Discrete Mathematics \textbf{140} (1995), 79--94.
\href{https://doi.org/10.1016/0012-365X(93)E0176-5}{doi:10.1016/0012-365X(93)E0176-5}.

\bibitem{Lacz12}
M.~Laczkovich,
\emph{Tilings of convex polygons with congruent triangles},
Discrete \& Computational Geometry \textbf{48} (2012), 330--372.
\href{https://doi.org/10.1007/s00454-012-9404-x}{doi:10.1007/s00454-012-9404-x}.

\bibitem{SWW91}
S.~L.~Snover, C.~Waiveris, and J.~K.~Williams,
\emph{Rep-tiling for triangles},
Discrete Mathematics \textbf{91} (1991), 193--200.
\href{https://doi.org/10.1016/0012-365X(91)90110-N}{doi:10.1016/0012-365X(91)90110-N}.

\bibitem{Soifer09}
A.~Soifer,
\emph{How Does One Cut a Triangle?},
2nd ed., Springer, New York, 2009; see pp.~47--50.

\bibitem{Zhang25}
Y.~X.~Zhang,
\emph{Tiling triangles with $2\pi/3$ angles},
arXiv:2512.22696v4 (2026).

\end{thebibliography}

\end{document}
